\section{Recall and the choice of parameters}
\subsection{Recovering the reference results}
We now have ways to avoid some scores, but how do we check whether we kept the right documents?
Let $G_K(q)$ be the reference top-$K$ IDs from a complete scan, and $A_K(q)$ the IDs returned
by the method we're testing. For $K\le N$, recall for query $q$ counts the reference IDs we recover:
\begin{equation}
R_K(q)=\frac{|G_K(q)\cap A_K(q)|}{K}.
\label{eq:recall}
\end{equation}
In our six-document example, $G_2=[3,4]$. If we return [3,0], we recover one of the two IDs,
so recall is 50\%; returning [3,4] gives 100\%. A large candidate set doesn't settle this
question by itself. It has to contain the reference IDs, and the final selection has to keep them.

We also need to be clear about which reference we're using. The main experiments compare
against the full binary-document score in Equation~\ref{eq:score}. If we instead used a
float-vector reference, we would also measure the loss from quantization. Human relevance
labels ask another question: whether the text is useful for the query. Recovering the binary
neighbors doesn't by itself establish either relevance or accurate language-model answers.

If we follow a reference document through the pipeline, we can lose it when its cluster isn't
opened, or when local search misses it inside an opened cluster. Fix a query and choose one
of its $K$ reference IDs uniformly. Let $V$ mean that we visit its cluster, and $L$ mean that
the local method returns it. Then
\begin{equation}
R_K(q)=\Pr(V\mid q)\Pr(L\mid V,q).
\label{eq:survival}
\end{equation}
This uses conditional probability; it doesn't assume the two events are independent. And if
we average over queries, we generally can't multiply the two separately averaged probabilities
instead. A's local scan has conditional survival one, but routing can still remove a result.
B and C share that routing limit, with any local work limits adding another way to miss it.

\subsection{Derive one useful distributional model}
To estimate recall mathematically, we need assumptions about both the document signs and the
query weights. Let's start with independent document signs, each equally likely to be 0 or 1.
Give $m$ strong query coordinates magnitude $a$, and the remaining $d-m$ coordinates magnitude
$\epsilon$. For this example, require
\begin{equation}
a>(d-m)\epsilon.
\label{eq:dominance}
\end{equation}
Here one strong mismatch costs more than all possible weak mismatches together. So a document
matching all the strong coordinates outranks every document that mismatches one, regardless
of what happens on the weak coordinates. We've deliberately constructed this condition;
Matryoshka training alone doesn't give us that guarantee.

If the signs are independent and balanced, a document matches those $m$ signs with probability
$p=2^{-m}$. We can then count how many matches occur among $N$ rows. Call that count $Z$;
it follows a binomial distribution:
\[
\Pr(Z=z)=\binom Nz p^z(1-p)^{N-z},\qquad \E[Z]=Np.
\]
Suppose we open only this ideal strong-sign group and score all its rows. If $Z\ge K$, the
whole top-$K$ is inside it. If $Z<K$, we recover $Z$ of those results. The recall is therefore
$\min(Z,K)/K$, and averaging over the possible counts gives
\begin{equation}
\E[R_K]=\frac1K\sum_{z=0}^{N}\min(z,K)\binom Nz p^z(1-p)^{N-z}
       =\frac1K\sum_{i=1}^{K}\Pr(Z\ge i).
\label{eq:binomial-recall}
\end{equation}
We can read the second form as adding the probabilities of recovering at least one result,
at least two, and so on through $K$. That lets us avoid a sum over every possible corpus count.

Consider $N=2$, $m=1$, and $K=1$. The ideal group is empty with probability $1/4$, contains
one document with probability $1/2$, and contains two with probability $1/4$. We can calculate
expected recall directly: $0(1/4)+1(1/2)+1(1/4)=3/4$. Notice that the average group size is
one, but some groups are empty. If we substituted that average as though it were a guaranteed
count, we would incorrectly predict 100\% recall.

\fig{recall-example}{Two independent documents and one required matching bit. Opening only the matching group succeeds with probability $3/4$. The empty-group outcome explains why the mean group size is insufficient.}{.84}

At $N=10^6$ and $m=13$, the ideal group averages about 122 rows. With $K=100$, some groups
can still be too small, so a complete B or C search may need to continue into other groups.
Equation~\ref{eq:binomial-recall} describes stopping after the ideal group; it doesn't describe
every search policy. If we increase a work budget, we have to consider which additional groups
we visit and the scores and sizes of those groups.

For ordinary queries, taking a fraction of the coordinates doesn't mean we keep the same
fraction of useful ranking information. Correlated bits, different query magnitudes, and
approximate routing all affect which neighbors survive. We can try a normal approximation
to the prefix and full scores, but then we need to calibrate it and check it independently.
In the earlier local Nomic check, recall errors reached about 9.6 percentage points, missing
the proposed two-point criterion. So that approximation cannot support a claim about real-data
recall at a billion rows here.

\subsection{Which parameters can we change?}
Some choices change what we build: the code width $d$, cluster count $C$, routing method,
and C's key coordinates and width $h$. Others can change from query to query: probes $P$,
B's leaf size and node budget, C's lookup and candidate limits, and the requested $K$. If
we change $d$, the reference ranking changes too, so we need to compare all methods again
using that same representation.

We can use the cost equation to narrow the settings worth trying. Suppose B follows one
promising path, each split halves the candidates, and the masks keep width $W=\ceil{N/w}$.
After $j$ splits, we expect $N2^{-j}$ rows to reach the scorer. If we temporarily leave out
changes in queue cost, this gives
\begin{equation}
T_B(j)\approx T_0+jWc_{\rm split}+Wc_{\rm leaf}+N2^{-j}c_{\rm score}.
\label{eq:depth}
\end{equation}
If we make one more split, we save about $N2^{-(j+1)}$ scores and pay for another full-width
split. It helps in this model when the saved scoring time exceeds $Wc_{\rm split}$ and the
additional queue cost. For example, with $N=6{,}400$ and $j=5$, another balanced cut saves
about 100 scores. We compare the time of those 100 scores with the time of a 100-word split;
the two counts alone aren't comparable.

Treating $j$ as continuous gives a starting point:
\[
\frac{dT_B}{dj}=Wc_{\rm split}-N\ln(2)2^{-j}c_{\rm score}=0,
\qquad
j^*=\log_2\!\left(\frac{N\ln(2)c_{\rm score}}{Wc_{\rm split}}\right).
\]
This gives us a place to start, and then we check the neighboring valid integer depths, the
score assumptions, and the required recall. The calculation assumes one useful balanced path;
it doesn't cover arbitrary branch exploration. Also, since $W$ grows with $N$, a larger
collection doesn't automatically call for a larger optimal depth. The dense masks themselves
still take more work as the collection grows.

The same distribution gives us a reason to try a different leaf size. At one million rows,
the 13-bit ideal group has mean about 122 and standard deviation about 11. If we stop at 128,
some slightly larger groups force another split on a weak coordinate. The rough margin
$122+3(11)\approx155$ suggested testing $L=160$. Since the preceding 12-bit group averages
about 244 rows, it will usually still split. We checked this suggestion experimentally;
the approximation alone doesn't guarantee that it is the right threshold.

The complete choice for method $a$ is
\begin{equation}
\min_{\theta\in\Theta_a}T_a(\theta)
\quad\text{subject to}\quad R_a(\theta)\ge r,\quad M_a(\theta)\le M_{\rm limit}.
\label{eq:choice}
\end{equation}
Here the vector $\theta$ holds the settings available to a method. A derivative is useful
when we know the relevant functions, but $N$ and RAM alone don't tell us cluster occupancy
or routing recall. We therefore measure promising settings, keep the feasible ones, and
compare the best \emph{tested} choices. That includes giving A enough routing options and
letting B stop splitting when a scan-like leaf is cheaper.